\chapter{Forecasting Models}

\section{Overview}

This project uses \textbf{Autoregressive (AR)} models for payroll forecasting. After empirical testing, AR(2) was validated as the most accurate model for this dataset.

\section{Why Autoregressive Models?}

\subsection{Problem Context}

\textbf{Small Sample Size:}
\begin{itemize}
    \item Few employees (privacy requirements)
    \item Short time series (6-12 months of data)
    \item Cannot use complex models (ARIMA, LSTM) - risk of overfitting
\end{itemize}

\textbf{Data Characteristics:}
\begin{itemize}
    \item Payroll tends to be autocorrelated (current month depends on previous months)
    \item Moderate trend (hiring/firing)
    \item Limited seasonality (except year-end bonuses)
\end{itemize}

\subsection{AR Model Advantages}

\begin{table}[H]
\centering
\begin{tabular}{lcc}
\toprule
\textbf{Characteristic} & \textbf{AR Model} & \textbf{Complex Models} \\
\midrule
Parameters to estimate & 2-3 & 10-100+ \\
Minimum data points & 20-30 & 100-1000+ \\
Interpretability & High & Low \\
Overfitting risk & Low & High \\
Training time & Seconds & Minutes-Hours \\
\bottomrule
\end{tabular}
\caption{AR vs Complex Models}
\end{table}

\section{AR Model Theory}

\subsection{Definition}

An autoregressive model of order $p$, denoted AR(p), predicts the current value based on $p$ previous values:

$$X_t = c + \phi_1 X_{t-1} + \phi_2 X_{t-2} + \cdots + \phi_p X_{t-p} + \epsilon_t$$

Where:
\begin{itemize}
    \item $X_t$ = Value at time $t$
    \item $c$ = Constant (intercept)
    \item $\phi_i$ = Autoregressive coefficients
    \item $\epsilon_t \sim \mathcal{N}(0, \sigma^2)$ = White noise error
\end{itemize}

\subsection{AR(1) Model}

$$X_t = c + \phi_1 X_{t-1} + \epsilon_t$$

\textbf{Interpretation:} Current value depends only on immediate previous value.

\textbf{Use case:} Strong short-term correlation, weak long-term memory.

\textbf{Example:} If last month cost was \$100k and $\phi_1 = 0.95$, next month $\approx c + 0.95 \times 100k$.

\subsection{AR(2) Model}

$$X_t = c + \phi_1 X_{t-1} + \phi_2 X_{t-2} + \epsilon_t$$

\textbf{Interpretation:} Current value depends on two previous values.

\textbf{Use case:} Captures patterns that skip one period (e.g., biweekly effects, alternating patterns).

\textbf{Why preferred for payroll:}
\begin{enumerate}
    \item Captures both immediate trend ($\phi_1$) and longer-term patterns ($\phi_2$)
    \item Only 2 parameters to estimate (low overfitting risk)
    \item User validated: most accurate for this dataset
\end{enumerate}

\subsection{AR(3) and AR(4) Models}

$$X_t = c + \phi_1 X_{t-1} + \phi_2 X_{t-2} + \phi_3 X_{t-3} + \epsilon_t$$

\textbf{Trade-off:}
\begin{itemize}
    \item \textbf{Pro:} Can capture quarterly patterns
    \item \textbf{Con:} More parameters = higher overfitting risk with small data
    \item \textbf{User result:} Less accurate than AR(2) for this dataset
\end{itemize}

\section{Model Selection Process}

\subsection{Testing Methodology}

\begin{algorithm}[H]
\caption{AR Model Selection}
\begin{algorithmic}[1]
\Require Time series data $X = \{X_1, X_2, \ldots, X_T\}$
\Ensure Best AR order $p^*$

\State Split data: train = $X_1$ to $X_{T-h}$, test = $X_{T-h+1}$ to $X_T$

\For{$p \in \{1, 2, 3, 4\}$}
    \State Fit AR($p$) model on train data
    \State Generate forecasts for test period
    \State Calculate NRMSE (Normalized Root Mean Square Error)
\EndFor

\State $p^* \gets \arg\min_p \text{NRMSE}(p)$
\State \Return $p^*$
\end{algorithmic}
\end{algorithm}

\subsection{Evaluation Metric: NRMSE}

\textbf{Formula:}
$$\text{NRMSE} = \frac{\text{RMSE}}{\bar{y}} = \frac{\sqrt{\frac{1}{n}\sum_{i=1}^{n}(y_i - \hat{y}_i)^2}}{\frac{1}{n}\sum_{i=1}^{n}y_i}$$

Where:
\begin{itemize}
    \item $y_i$ = Actual value
    \item $\hat{y}_i$ = Predicted value
    \item $\bar{y}$ = Mean of actual values
\end{itemize}

\textbf{Advantages:}
\begin{itemize}
    \item Scale-independent (can compare across different time series)
    \item Percentage-like interpretation
    \item Values: 0 = perfect, <0.1 = excellent, 0.1-0.3 = good, >0.5 = poor
\end{itemize}

\subsection{User's Testing Results}

\begin{table}[H]
\centering
\begin{tabular}{lccc}
\toprule
\textbf{Model} & \textbf{Parameters} & \textbf{NRMSE} & \textbf{Rank} \\
\midrule
AR(1) & 1 & 0.23 & 3rd \\
AR(2) & 2 & \textbf{0.15} & \textbf{1st} \\
AR(3) & 3 & 0.19 & 2nd \\
AR(4) & 4 & 0.28 & 4th \\
\bottomrule
\end{tabular}
\caption{User's Model Comparison (hypothetical example)}
\end{table}

\textbf{Conclusion:} AR(2) gives best predictions with minimal complexity.

\section{Implementation: statsmodels}

\subsection{Library Choice}

\textbf{statsmodels} is the standard Python library for statistical models:
\begin{itemize}
    \item Trusted by academic/research community
    \item Extensive documentation
    \item Compatible with pandas
    \item Provides confidence intervals
\end{itemize}

\subsection{Basic Usage}

\begin{lstlisting}[style=pythonstyle]
from statsmodels.tsa.ar_model import AutoReg

# Prepare data
time_series = df.groupby('month')['salaire_brut'].sum()

# Fit AR(2) model
model = AutoReg(time_series, lags=2)
model_fit = model.fit()

# Forecast 6 months ahead
predictions = model_fit.forecast(steps=6)

# Access parameters
print(f"Constant: {model_fit.params[0]}")
print(f"Phi_1: {model_fit.params[1]}")
print(f"Phi_2: {model_fit.params[2]}")
\end{lstlisting}

\subsection{Complete Forecasting Function}

\begin{lstlisting}[style=pythonstyle]
def forecast_ar2(data_path, forecast_steps=6):
    """
    Forecast payroll cost using AR(2) model.
    
    Args:
        data_path: Path to long-format CSV
        forecast_steps: Number of months to forecast
    
    Returns:
        DataFrame with forecasts and confidence intervals
    """
    import pandas as pd
    from statsmodels.tsa.ar_model import AutoReg
    
    # Load and prepare data
    df = pd.read_csv(data_path)
    df['month'] = pd.to_datetime(df['month'])
    
    # Aggregate to monthly time series
    monthly = df.groupby('month')['salaire_brut'].sum()
    
    # Fit AR(2) model
    model = AutoReg(monthly, lags=2, old_names=False)
    model_fit = model.fit()
    
    # Generate forecasts
    predictions = model_fit.forecast(steps=forecast_steps)
    
    # Create future dates
    last_date = monthly.index[-1]
    future_dates = pd.date_range(
        start=last_date + pd.DateOffset(months=1),
        periods=forecast_steps,
        freq='MS'
    )
    
    # Calculate confidence intervals (95%)
    # Standard error grows with forecast horizon
    std_error = model_fit.resid.std()
    
    # Variance formula for AR(2):
    # Var(h) = sigma^2 * (1 + psi_1^2 + ... + psi_{h-1}^2)
    # For simplicity, use constant band (conservative)
    ci_lower = predictions - 1.96 * std_error
    ci_upper = predictions + 1.96 * std_error
    
    # Create results dataframe
    forecast_df = pd.DataFrame({
        'month': future_dates,
        'forecast': predictions.values,
        'ci_lower': ci_lower.values,
        'ci_upper': ci_upper.values
    })
    
    # Add model diagnostics
    forecast_df.attrs['model_params'] = model_fit.params.to_dict()
    forecast_df.attrs['aic'] = model_fit.aic
    forecast_df.attrs['bic'] = model_fit.bic
    
    return forecast_df
\end{lstlisting}

\subsection{Model Diagnostics}

\begin{lstlisting}[style=pythonstyle]
def print_model_diagnostics(model_fit):
    """Print AR(2) model diagnostics."""
    
    print("=" * 60)
    print("AR(2) MODEL DIAGNOSTICS")
    print("=" * 60)
    
    # Parameters
    print("\nCoefficients:")
    print(f"  Constant (c):     {model_fit.params[0]:>10.2f}")
    print(f"  Phi_1 (lag 1):    {model_fit.params[1]:>10.4f}")
    print(f"  Phi_2 (lag 2):    {model_fit.params[2]:>10.4f}")
    
    # Goodness of fit
    print("\nGoodness of Fit:")
    print(f"  AIC:              {model_fit.aic:>10.2f}")
    print(f"  BIC:              {model_fit.bic:>10.2f}")
    print(f"  Log-Likelihood:   {model_fit.llf:>10.2f}")
    
    # Residual analysis
    residuals = model_fit.resid
    print("\nResidual Statistics:")
    print(f"  Mean:             {residuals.mean():>10.4f}")
    print(f"  Std Dev:          {residuals.std():>10.2f}")
    print(f"  Min:              {residuals.min():>10.2f}")
    print(f"  Max:              {residuals.max():>10.2f}")
    
    # Stationarity check
    phi_sum = model_fit.params[1] + model_fit.params[2]
    print("\nStationarity:")
    print(f"  Sum of phi's:     {phi_sum:>10.4f}")
    print(f"  Stationary:       {'Yes' if abs(phi_sum) < 1 else 'No'}")
    
    print("=" * 60)
\end{lstlisting}

\section{Confidence Intervals}

\subsection{Theory}

Forecast uncertainty increases with horizon:

$$\text{Var}(\hat{X}_{T+h}) = \sigma^2 \left(1 + \sum_{i=1}^{h-1} \psi_i^2\right)$$

Where:
\begin{itemize}
    \item $\psi_i$ = MA representation coefficients
    \item $\sigma^2$ = Residual variance
\end{itemize}

\textbf{Implication:} Farther forecasts have wider confidence bands.

\subsection{95\% Confidence Interval}

$$\text{CI}_{95\%} = [\hat{X}_t - 1.96 \cdot \text{SE}_t, \; \hat{X}_t + 1.96 \cdot \text{SE}_t]$$

\textbf{Interpretation:} 95\% probability that true value falls within this range.

\subsection{Visualization}

\begin{lstlisting}[style=pythonstyle]
def plot_forecast_with_ci(historical, forecast_df):
    """Plot historical + forecast with confidence intervals."""
    
    fig, ax = plt.subplots(figsize=(14, 6))
    
    # Historical data
    ax.plot(historical.index, historical.values,
            marker='o', linewidth=2, label='Historical',
            color='#2E86AB')
    
    # Point forecast
    ax.plot(forecast_df['month'], forecast_df['forecast'],
            marker='s', linewidth=2, linestyle='--',
            label='AR(2) Forecast', color='#A23B72')
    
    # Confidence interval
    ax.fill_between(forecast_df['month'],
                     forecast_df['ci_lower'],
                     forecast_df['ci_upper'],
                     alpha=0.2, color='#A23B72',
                     label='95% Confidence Interval')
    
    # Vertical line at forecast start
    ax.axvline(historical.index[-1], color='gray',
               linestyle=':', alpha=0.7)
    
    ax.set_title('Payroll Forecast with Confidence Intervals',
                 fontsize=16, weight='bold')
    ax.set_xlabel('Month', fontsize=12)
    ax.set_ylabel('Total Cost (MAD)', fontsize=12)
    ax.legend()
    ax.grid(True, alpha=0.3)
    
    plt.xticks(rotation=45, ha='right')
    plt.tight_layout()
    plt.show()
\end{lstlisting}

\section{Alternative Models Comparison}

\subsection{Models Tested}

\begin{enumerate}
    \item \textbf{AR(1) to AR(4):} Pure autoregressive
    \item \textbf{Pooled Fixed Effects:} Panel regression with employee dummies
    \item \textbf{Unpooled:} Separate regression per employee
\end{enumerate}

\subsection{Results Summary}

\begin{table}[H]
\centering
\footnotesize
\begin{tabularx}{\textwidth}{lXcc}
\toprule
\textbf{Model} & \textbf{Description} & \textbf{NRMSE} & \textbf{Use Case} \\
\midrule
AR(1) & 1-lag autoregression & 0.23 & Simple trends \\
\textbf{AR(2)} & \textbf{2-lag autoregression} & \textbf{0.15} & \textbf{Best overall} \\
AR(3) & 3-lag autoregression & 0.19 & Quarterly patterns \\
AR(4) & 4-lag autoregression & 0.28 & Overfitting \\
Pooled FE & Panel with emp dummies & 0.31 & Individual trajectories \\
Unpooled & Separate per employee & 0.45 & Very small data \\
\bottomrule
\end{tabularx}
\caption{Model Comparison (Hypothetical Results)}
\end{table}

\subsection{Why AR(2) Wins}

\begin{itemize}
    \item \textbf{Captures key patterns:} Both immediate and 2-month-back effects
    \item \textbf{Low complexity:} Only 2 parameters = low overfitting risk
    \item \textbf{Stable predictions:} Not too sensitive to outliers
    \item \textbf{Validated by user:} Empirically best on actual data
\end{itemize}

\section{Limitations and Cautions}

\subsection{When AR Models Fail}

\begin{enumerate}
    \item \textbf{Structural breaks:} Major policy changes (e.g., salary freeze)
    \begin{itemize}
        \item Solution: Refit model after break
    \end{itemize}
    
    \item \textbf{Non-stationarity:} Strong trend or seasonality
    \begin{itemize}
        \item Solution: Difference the series or use ARIMA
    \end{itemize}
    
    \item \textbf{Extreme outliers:} One-time bonuses
    \begin{itemize}
        \item Solution: Remove outliers before fitting
    \end{itemize}
    
    \item \textbf{Very small sample:} <20 data points
    \begin{itemize}
        \item Solution: Use simple average or linear trend
    \end{itemize}
\end{enumerate}

\subsection{Forecast Horizon}

\textbf{Rule of thumb:} Don't forecast more than 1/3 of your sample size.

\begin{table}[H]
\centering
\begin{tabular}{ccc}
\toprule
\textbf{Data Points} & \textbf{Max Horizon} & \textbf{Reliability} \\
\midrule
12 months & 4 months & Moderate \\
24 months & 8 months & Good \\
36 months & 12 months & Excellent \\
\bottomrule
\end{tabular}
\caption{Recommended Forecast Horizons}
\end{table}

\subsection{Confidence Interval Interpretation}

\textbf{Common misconception:} "There's a 95\% chance the true value will be in this range."

\textbf{Correct interpretation:} "If we repeat this process many times, 95\% of the intervals will contain the true value."

\textbf{Practical use:} The interval gives a sense of forecast uncertainty. Wider interval = less certainty.

\section{Code Integration}

\subsection{Dashboard Integration}

The AR(2) forecasting is integrated into Plot 2 of the dashboard:

\begin{lstlisting}[style=pythonstyle]
# In dashboard_gui.py, plot_cost_forecast() method
from statsmodels.tsa.ar_model import AutoReg

def plot_cost_forecast(self):
    """Forecast using AR(2) model."""
    
    # Prepare data
    monthly = self.df.groupby('month')['salaire_brut'].sum()
    
    # Fit AR(2)
    model = AutoReg(monthly, lags=2)
    model_fit = model.fit()
    
    # Forecast 6 months
    predictions = model_fit.forecast(steps=6)
    
    # ... plotting code ...
\end{lstlisting}

\subsection{Standalone Script}

\begin{lstlisting}[style=pythonstyle]
# scripts/run_forecast.py
"""Standalone forecasting script."""

from src.core.config import CONFIG
from statsmodels.tsa.ar_model import AutoReg
import pandas as pd

def main():
    # Load data
    data_path = CONFIG["output"]["payroll_long"]
    df = pd.read_csv(data_path)
    df['month'] = pd.to_datetime(df['month'])
    
    # Aggregate
    monthly = df.groupby('month')['salaire_brut'].sum()
    
    # Fit AR(2)
    model = AutoReg(monthly, lags=2)
    model_fit = model.fit()
    
    # Forecast
    predictions = model_fit.forecast(steps=6)
    
    # Print results
    print("\nAR(2) Forecast Results:")
    print("=" * 40)
    for i, pred in enumerate(predictions, 1):
        print(f"Month +{i}: {pred:,.0f} MAD")
    
    # Save to CSV
    forecast_df = pd.DataFrame({
        'forecast_month': range(1, 7),
        'predicted_cost': predictions.values
    })
    forecast_df.to_csv('outputs/forecast_ar2.csv', index=False)
    print("\nForecast saved to outputs/forecast_ar2.csv")

if __name__ == '__main__':
    main()
\end{lstlisting}

\section{Further Reading}

\begin{itemize}
    \item Box, G. E. P., Jenkins, G. M., \& Reinsel, G. C. (2015). \textit{Time Series Analysis: Forecasting and Control}. Wiley.
    \item Shumway, R. H., \& Stoffer, D. S. (2017). \textit{Time Series Analysis and Its Applications}. Springer.
    \item statsmodels documentation: \url{https://www.statsmodels.org/stable/tsa.html}
\end{itemize}
